\section{Additional Serving Results}
\label{app:serving-more}

This appendix collects serving measurements that Section~\ref{sec:cost} does not show. They come from the same two
setups: one $4\times\text{GB200}$ tray (TP4/EP4) and one $8\times\text{H200}$ node (TP8/EP8), both running SGLang
\texttt{da64c5cbb} with our plugin. Runs marked \emph{dummy weights} use random weights with standard deviation 0.02,
which route tokens to experts almost uniformly. In an A/B test, dummy weights decoded 15--36\% slower than the real
ones, while the ratios between models moved by 2--4\%. All other runs use the original (real) \vfour-Flash weights, unchanged. In Figure~\ref{fig:serving-main}, panel (a) uses dummy weights and
panels (b) and (c) use the real weights. Every
comparison is at equal rounds $T$, and the names are those of Section~\ref{sec:cost}. Across layers, a model either
shares global KV as \ced does or keeps \emph{per-layer KV} (every layer \Lr{2}{39} caches its own global KV and runs
over every prompt token in prefill). Across rounds, it keeps \emph{last-round KV} or \emph{per-round KV}. The
\emph{baseline} keeps per-layer and per-round KV; \emph{\ourmethod} keeps neither. The \emph{decoder loop} loops
\Lr{21}{38} of \ced instead of the encoder. Per-round KV is always emulated by shrinking the KV pool, so a model and its
per-round variant run identical compute.

\begin{figure*}[t]
\centering
\safegraphic[width=\textwidth]{figures/fig_serving_more.pdf}
\caption{\textbf{Further serving measurements.} \emph{(a)}~As Figure~\ref{fig:serving-main}(a), for ShareGPT on
one GB200 tray (dummy weights): output tokens/s per GPU against the p90 per-user decode speed
($1000/\text{p90 TPOT}$), one point per concurrency, preemption-free points only. \emph{(b, c)}~Mooncake tool-agent trace, closed loop, $T=8$, real
weights (the runs of Figure~\ref{fig:p1-pareto}c): \emph{(b)}~prefix-cache hit rate over the second half of each run;
\emph{(c)}~share of requests with TTFT $\le 10$\,s.}
\label{fig:serving-more}
\end{figure*}

\paragraph{The gain holds at interactive speeds.}
The $3.1\times$--$3.6\times$ of Section~\ref{sec:cost} compares peak throughput, which \ourmethod reaches at $c=512$.
Figures~\ref{fig:serving-main}(a, b) and~\ref{fig:serving-more}(a) plot throughput against the p90 per-user speed
instead. Reading the curves by linear interpolation between neighbouring points, at 30 and 50 tokens/s per user \ourmethod serves $1.39\times$--$1.67\times$ the throughput of
the baseline on 8K/1K and $1.42\times$--$1.61\times$ on ShareGPT. No KV pool is full at these speeds, so this part of
the gain comes from cross-layer sharing, not from capacity: the baseline runs every layer on every prompt token
(166 instead of 147 layer passes at $T=8$) and keeps a KV cache and an indexer in all 38 layers. A single request
shows the same: its time to first token (TTFT) on 8K/1K is 0.98\,s instead of 1.47\,s at $T=4$ and 1.68\,s instead of
2.65\,s at $T=8$, and its time per output token is 10\% lower. On H200 with real weights the gap is larger. At equal
concurrency, \ourmethod serves $1.48\times$--$1.64\times$ the throughput of the baseline on 8K/1K ($c=1$--64) and
$1.47\times$--$1.73\times$ on ShareGPT ($c=1$--256); its single-request TTFT is 1.95\,s instead of 3.23\,s ($T=4$)
and 3.56\,s instead of 5.61\,s ($T=8$). On 8K/1K the H200 baseline preempts at $c=112$ ($T=4$) and fills 96\% of its
pool at $c=64$ ($T=8$), while \ourmethod runs $c=256$ at 16\% KV use. On H200 ShareGPT the baseline's pool never
exceeds 28\%, so the $1.72\times$--$1.73\times$ there at $c=256$ comes from cross-layer sharing alone.

\begin{table*}[t]
\caption{\textbf{Cross-layer and cross-round sharing, measured separately} ($4\times\text{GB200}$, dummy weights).
Left: peak preemption-free output tokens/s per GPU on 8K/1K, with the largest preemption-free concurrency in
parentheses. Right: decode tokens/s per GPU at the largest batch of 128K-token requests that fits the KV pool, with
that batch in parentheses. Without a loop, \ced reaches 887.8 and 780.3 and the per-layer-KV model 371.7 and 116.7.
The ratios compare rows at equal $T$.}
\label{tab:serving-2x2}
\vskip 0.06in
\centering
\footnotesize
\setlength{\tabcolsep}{5pt}
\begin{tabular}{@{}llcccc@{}}
\toprule
& & \multicolumn{2}{c}{8K in / 1K out, peak} & \multicolumn{2}{c}{128K decode at the KV limit} \\
\cmidrule(lr){3-4}\cmidrule(lr){5-6}
KV across layers & KV across rounds & $T=4$ & $T=8$ & $T=4$ & $T=8$ \\
\midrule
\ced (\ourmethod, ours) & last round & 391.8 ($c$\,512) & 185.1 ($c$\,512) & 347.4 (38) & 199.2 (38) \\
\ced & per round & 308.4 ($c$\,128) & 142.5 ($c$\,64) & 186.1 (13) & 69.4 (7) \\
per layer & last round & 161.6 ($c$\,32) & 92.4 ($c$\,32) & 51.7 (3) & 29.6 (3) \\
per layer (baseline) & per round & 127.3 ($c$\,16) & 50.9 ($c$\,8) & 21.1 (1) & does not fit (0) \\
\midrule
\multicolumn{2}{@{}l}{Cross-layer sharing, both last-round KV (\ced $\div$ per layer)} & $2.43\times$ & $2.00\times$ & $6.72\times$ & $6.73\times$ \\
\multicolumn{2}{@{}l}{Cross-layer sharing, both per-round KV} & $2.42\times$ & $2.80\times$ & $8.82\times$ & --- \\
\multicolumn{2}{@{}l}{Cross-round sharing in \ced (last round $\div$ per round)} & $1.27\times$ & $1.30\times$ & $1.87\times$ & $2.87\times$ \\
\multicolumn{2}{@{}l}{Cross-round sharing with per-layer KV} & $1.27\times$ & $1.82\times$ & $2.45\times$ & --- \\
\multicolumn{2}{@{}l}{Both (\ourmethod $\div$ baseline)} & $\mathbf{3.08\times}$ & $\mathbf{3.64\times}$ & $\mathbf{16.5\times}$ & 199.2 vs.\ none \\
\bottomrule
\end{tabular}
\end{table*}

\paragraph{Both kinds of sharing contribute, and they multiply.}
Table~\ref{tab:serving-2x2} measures all four combinations of cross-layer sharing (\ced or per-layer KV) and
cross-round sharing (last-round or per-round KV). With last-round KV in both, cross-layer sharing alone gives
$2.0\times$--$2.4\times$ the peak 8K/1K throughput and $6.7\times$ the 128K decode throughput, because \ced fits 38
requests of 128K tokens where per-layer KV fits 3. Cross-round sharing alone gives $1.27\times$--$1.30\times$ in \ced
and $1.27\times$--$1.82\times$ with per-layer KV on 8K/1K, and $1.87\times$--$2.87\times$ in \ced at 128K. The 8K/1K
factors multiply to the headline exactly ($1.27\times 2.43=3.08$, $1.82\times 2.00=3.64$). Cross-round sharing thus
also helps a model with per-layer KV: the benefit of stochastic-depth training is not specific to \ced. At equal batch,
a model and its per-round variant decode within 3\% of each other (one point differs by 6\%), as the emulation intends,
so these gains come from the larger batch that fits. For the same reason, cross-round sharing changes nothing in \ced
at 30 or 50 tokens/s per user (factor 1.00), where its pool is at most a third full. It pays off once the pool fills:
at long context, on agent traces and under heavy load.

\begin{table*}[t]
\caption{\textbf{Prefill cost of running every layer over the prompt.} Top: ratio of per-request prefill time for
uncached prompts (batch 1, real weights); both models keep last-round KV. \emph{Full replay} runs the \ced decoder over
every prompt token instead of the last 128. Bottom: ratio of new prompt tokens prefilled per second behind one cached
prefix ($4\times\text{GB200}$, dummy weights, 4K new tokens per request). Values above 1 mean the \ced variant is
faster.}
\label{tab:serving-prefill}
\vskip 0.06in
\centering
\footnotesize
\setlength{\tabcolsep}{5pt}
\begin{tabular}{@{}llccccc@{}}
\toprule
& & \multicolumn{5}{c}{Prompt length (tokens)} \\
\cmidrule(lr){3-7}
& Ratio & 512 & 2K & 4K & 16K & 64K \\
\midrule
H200 & \ced, no loop: full replay $\div$ 128-token tail & 1.32 & 1.45 & 1.50 & 1.60 & 1.75 \\
H200 & Per-layer KV $\div$ \ced, encoder loop $T=8$ & 1.87 & 1.45 & 1.44 & 1.69 & 2.43 \\
GB200 & Per-layer KV $\div$ \ced, encoder loop $T=8$ & 1.65 & 1.48 & 1.41 & 1.40 & 1.48 \\
\midrule
& & \multicolumn{5}{c}{Context length / prefill chunk (tokens)} \\
\cmidrule(lr){3-7}
& & 64K / 4K & 64K / 8K & 64K / 16K & 128K / 4K & 128K / 8K \\
\midrule
GB200 & \ourmethod $\div$ baseline, $T=4$ & 1.60 & 1.64 & 1.70 & 1.65 & 1.83 \\
\bottomrule
\end{tabular}
\end{table*}

\paragraph{Prefill: skipping the decoder survives looping.}
A \ced prompt runs the encoder and a 128-token decoder tail, and an encoder loop keeps that property
(Table~\ref{tab:serving-prefill}). Even without a loop, running the decoder over the whole prompt costs $1.32\times$ at
512 tokens and $1.75\times$ at 64K tokens on H200, and $1.30\times$ at 8K tokens on GB200. With eight rounds, the
model with per-layer KV takes $1.40\times$--$1.65\times$ (GB200) and $1.44\times$--$2.43\times$ (H200) as long as the
\ced encoder loop to prefill the same prompt. Behind a cached 60K--124K-token prefix, \ourmethod at $T=4$ prefills
$1.60\times$--$1.83\times$ as many new tokens per second as the baseline at the same chunk size, and $1.56\times$ at
$T=8$ (64K context, 4K chunk). The encoder loop's own prefill grows more slowly than its round count
(Table~\ref{tab:unit-costs-app}).

\begin{table*}[t]
\caption{\textbf{Requests that fit the KV pool at long context.} $4\times\text{GB200}$: one tray with the flags of
Figure~\ref{fig:p1-pareto}, dummy weights; requests of $L+256$ tokens; the last column is decode tokens/s per GPU at
the largest 128K batch that fits. The \ced limit of 604 at 8K is computed from the pool (the sweep stopped at 256).
$8\times\text{H200}$: one node, real weights, requests that fit at full length. The H200 per-round pool also charges
each round's sliding-window rings; the GB200 one does not.}
\label{tab:serving-longctx}
\vskip 0.06in
\centering
\footnotesize
\setlength{\tabcolsep}{5pt}
\begin{tabular}{@{}llcccccc@{}}
\toprule
& & \multicolumn{3}{c}{$4\times\text{GB200}$} & \multicolumn{3}{c}{$8\times\text{H200}$} \\
\cmidrule(lr){3-5}\cmidrule(lr){6-8}
Model & KV across rounds & 8K & 128K & tok/s/GPU at 128K & 120K & 524K & 1M \\
\midrule
\ced, no loop & --- & 604 & 38 & 780.3 & 195 & 44 & 23 \\
\ced encoder loop, $T=4$ & last round & 604 & 38 & 347.4 & --- & --- & --- \\
\ced encoder loop, $T=8$ & last round & 604 & 38 & 199.2 & 195 & 44 & 23 \\
\ced decoder loop, $T=8$ & last round & --- & --- & --- & 195 & 44 & 23 \\
\ced encoder loop, $T=4$ & per round & 215 & 13 & 186.1 & --- & --- & --- \\
\ced encoder loop, $T=8$ & per round & 116 & 7 & 69.4 & 23 & 5 & 2 \\
\midrule
Per-layer KV, no loop & --- & 49 & 3 & 116.7 & 16 & 3 & 1 \\
Per-layer KV, $T=4$ & last round & 49 & 3 & 51.7 & --- & --- & --- \\
Per-layer KV, $T=8$ & last round & 49 & 3 & 29.6 & 16 & 3 & 1 \\
Per-layer KV, $T=4$ & per round & 25 & 1 & 21.1 & --- & --- & --- \\
Per-layer KV, $T=8$ & per round & 15 & 0 & --- & --- & --- & --- \\
\bottomrule
\end{tabular}
\end{table*}

\paragraph{Cost per request on four workloads.}
Figure~\ref{fig:serving-main}(c) runs four closed-loop workloads on one GB200 tray with real weights at $T=8$, one run
each:
synthetic reasoning (2K-token prompts, 4K output tokens, 32 concurrent requests), the Mooncake conversation and
tool-agent traces~\citep{mooncake} (64 concurrent requests, starting at minute 10), and synthetic long documents
(100K-token prompts, 1K output tokens, 8 concurrent requests). \ourmethod costs $3.7\times$--$4.7\times$ the GPU time
per request of an unlooped \ced. That is $0.64\times$--$0.93\times$ the cost of per-layer KV with last-round KV, and
$0.36\times$--$0.93\times$ the cost of the baseline. The largest saving is on long documents, where the baseline's pool
holds one 101K-token request at a time. The cost model, evaluated for the same runs, matches the two \ced loops
(measured over predicted 0.78--1.04). It underpredicts per-layer KV with last-round KV by 9--44\% on the three workloads
where that model's KV pool fills (95--100\%), an effect that a model of FLOPs and bytes does not capture. The encoder
loop costs the same as the decoder loop on reasoning ($1.00\times$) and $1.37\times$--$1.64\times$ as much on the three
prompt-heavy workloads, where its rounds run over every prompt token. These runs decode 8--64 requests at a time on one tray. With larger decode batches, as in
expert-parallel serving across many GPUs, each output token costs less and prompt tokens take a larger share of a
request, so the premium grows; running prompts at fewer rounds ($T_p<T_d$, Appendix~\ref{app:split}) lowers it.

\paragraph{Cross-round sharing leaves prefix caching as it is without a loop.}
Figure~\ref{fig:serving-more}(b, c) replays the Mooncake tool-agent trace in closed loop
with 64--512 concurrent sessions, real weights and $T=8$ (the runs of Figure~\ref{fig:p1-pareto}c); each point replays
600--1{,}536 complete requests from the trace. With last-round KV, both the encoder and
the decoder loop hit the prefix cache within 0.8 points as often as the unlooped \ced at every concurrency, so looping
does not change which session histories stay resident. With per-round KV, the 1.5M-token pool is full from $c=128$,
and the hit rate is 4.3, 6.4, 8.1 and 11.6 points lower than with last-round KV at $c=64$, 128, 256 and 512. The
share of requests with TTFT $\le 10$\,s then falls to 4.8\% at $c=256$ (58.3\% with last-round KV) and 1.4\% at
$c=512$ (46.0\%). A synthetic multi-turn agent workload on H200 (80 sessions of 4 turns, 32K--40K-token contexts,
0.2 requests/s) shows the same direction: per-round KV lowers the hit rate of the encoder loop from 72.3\% to 63.8\%
and its completed rate from 0.199 to 0.171 requests/s. Cross-layer sharing matters even more on this workload. In the
same runs, per-layer KV with last-round KV holds 2.0M tokens against \ced's 23.4M, so the session histories (about
3.2M tokens) do not stay resident: its hit rate is 3.4\% and it completes 0.039 requests/s, $5.1\times$ fewer than
\ourmethod.

\paragraph{Encoder and decoder loops decode at the same cost.}
Decode runs every layer under either loop placement. In a pure-decode sweep with real weights on GB200 that fills the KV
pool (8K, 16K and 32K-token contexts, 256 output tokens per request), the encoder loop and the decoder loop at $T=8$
reach 1{,}451 / 1{,}118 / 818 and 1{,}460 / 1{,}112 / 828 output tokens/s per GPU (the 8K values are capped at 768
requests). With last-round KV neither loop adds global KV, so both hold the same pool, and both take 33.1\,ms per decode
step at 8K context and batch 32.

\paragraph{Long and very long contexts.}
With last-round KV, a looped \ced fits as many requests as the unlooped \ced at every $T$
(Table~\ref{tab:serving-longctx}): 38 requests of 128K tokens on one GB200 tray, and 44 of 524K tokens and 23 of 1M
tokens on one H200 node, where the unlooped model with per-layer KV fits 3 and 1. At 524K tokens on H200, last-round KV
decodes 195 tokens/s per node at a batch of 32. Per-round KV fits only 5 such requests and peaks at 55.7 tokens/s at a
batch of 4, $3.5\times$ less. For every model, the long-context decode measurements use KV that is allocated without
being computed, so that they measure decode alone; on H200, decoding over such KV was within 3\% of decoding after a
real prefill.
